% ======================================================================
\section{Prescriptions and distributions}\label{app:distributions}
% ======================================================================

We collect here the conventions used for the singular longitudinal and transverse integrals.

\paragraph{Principal values.} Energy denominators that vanish inside the region of integration, such as $p^+\v k^2-k^+\v p^2$ in momentum space or $p^+(\v x-\v u)^2-k^+(\v x-\v z)^2$ in coordinate space, arise from transitions between states of equal free energy and are defined as Cauchy principal values, Eq.~\eqref{eq:PV}. Correspondingly, logarithms that arise from integrals over such denominators are logarithms of absolute values, $\log|\cdot|$. All finite terms of this paper refer to this prescription (Section~\ref{sec:fin_conv}).

\paragraph{Dimensional regularisation.} Transverse loop integrals are computed in $d=2-\epsilon$ dimensions, $\int\frac{d^2p}{(2\pi)^2}\to\mu^{\epsilon}\int\frac{d^dp}{(2\pi)^d}$, in the $\MS$ scheme with scale $\mu_{\MS}$. Scaleless integrals vanish, $\int d^dp/\v p^2=0$. Coordinate-space logarithms are written with the combination~\eqref{eq:c0}.

\paragraph{Plus prescriptions in $p^+$.} For a function $f(p^+)$ regular at $p^+=0$,
\begin{align}\label{eq:plusp}
\int_0^{P}dp^+\Big[\frac{1}{p^+}\Big]_+f(p^+)\equiv\int_0^{P}dp^+\,\frac{f(p^+)-f(0)}{p^+}\,.
\end{align}
The subtracted term, integrated from $\Lambda$, is $f(0)\log(P/\Lambda)$; when $f(0)$ is the JIMWLK (target-side) structure it is part of the JIMWLK term.

\paragraph{Plus prescriptions in the momentum fraction.} With $\xi=k^+/(p^++k^+)$ the momentum fraction of the measured gluon in its parent, and for a function $F(\xi)$,
\begin{align}\label{eq:plusxi}
\int_0^1d\xi\,\frac{\xi}{(1-\xi)_+}F(\xi)&\equiv\int_0^1d\xi\,\frac{\xi F(\xi)-F(1)}{1-\xi}\,,\notag\\
\int_0^1d\xi\,\Big[\frac{1-\xi}{\xi}\Big]_+F(\xi)&\equiv\int_0^1d\xi\,\frac{1-\xi}{\xi}\big[F(\xi)-F(0)\big]\,.
\end{align}
The first is the standard prescription at $\xi\to1$ (integrated gluon soft); the second subtracts the pole at $\xi\to0$ (measured gluon soft relative to its parent), which belongs to the projectile (BFKL) evolution. The combination of the two subtractions coincides with the ``double plus'' prescription of~\cite{Kovner:2023vsy}: with $N(\xi)=\xi^2+(1-\xi)^2+\xi^2(1-\xi)^2$, $N(0)=N(1)=1$,
\begin{align}
\int_0^1\frac{d\xi}{\xi_+(1-\xi)_+}\,N(\xi)\,F(\xi)&\equiv\int_0^1d\xi\Big[\frac{N(\xi)F(\xi)}{\xi(1-\xi)}-\frac{F(0)}{\xi}-\frac{F(1)}{1-\xi}\Big]\,,\notag\\
\frac{N(\xi)}{\xi(1-\xi)}&=\xi(1-\xi)+\frac{\xi}{1-\xi}+\frac{1-\xi}{\xi}\,,
\end{align}
and $\int_0^1\frac{d\xi}{\xi_+(1-\xi)_+}N(\xi)=-\frac{11}{6}$.

\paragraph{Integrals over $p^+$ up to $\vee\to\infty$.} Some terms have no dependence on $p^+$ other than an oscillating phase. Their integral converges for $\vee\to\infty$ in the sense of distributions,
\begin{align}\label{eq:distrib}
\int_0^\infty dp^+\,e^{-iap^+}=\pi\,\delta(a)-i\,\mathcal P\frac1a\,.
\end{align}

\paragraph{Recoil function.} When only the recoil phase of the measured gluon depends on the longitudinal momentum, the plus-prescription integral of the rows with $\mf B_3$ and $\mf C_1$ is the function (Appendix~\ref{app:closedforms})
\begin{align}\label{eq:Erecoil}
\mathcal E(s)=\int_0^1\frac{dt}{t}\big(e^{-ist}-1\big)\,,
\end{align}
which vanishes at $s=0$ and grows logarithmically for $|s|\to\infty$, $\mathcal E(s)=-\log|s|-\gamma_E-\frac{i\pi}{2}{\rm sign}(s)+O(1/s)$.

\subsection{Remarks on the scheme}\label{app:scheme}

This subsection collects three remarks on the scheme of Section~\ref{sec:dglap_rc}, which are not needed to state the result.

\paragraph{Polarisations in two dimensions.} The result is given in conventional dimensional regularisation, with $d_\perp=2-\epsilon$ gluon polarisations where a polarisation sum multiplies a pole. With two-dimensional polarisations and $d$-dimensional momenta (a dimensional-reduction-like scheme) the reduction $\tilde p^l\tilde p^i\to\delta^{li}\tilde p^2/d$ gives a factor $2/d$ instead of $d_\perp/d=1$: the constant of Row~I.II would be $\frac{64}{9}-\frac{\pi^2}{3}$ (a shift $-\frac{\alpha_sN_c}{6\pi}\times$LO), that of Row~I.V would change by the opposite amount, as it must by Eq.~\eqref{eq:dg_SEid} below, Row~II.IV by $-\mathcal C\int_0^1d\xi\,2N_c\,\xi(1-\xi)\,\Pi(\xi)$, and the coupling and the fragmentation function would have to be converted, $\alpha_s^{\rm DR}=\alpha_s^{\MS}\big(1+\frac{\alpha_sN_c}{12\pi}\big)$ and $D^{\rm DR}=D^{\MS}\otimes\big[\delta(1-\xi)+\frac{\alpha_s}{2\pi}\big(2N_c\xi(1-\xi)-\frac{N_c}{6}\delta(1-\xi)\big)\big]$; the coefficient $\frac{N_c}{6}$ of the $\delta$ term is fixed independently by the momentum sum rule and equals the coupling conversion. With these conversions the result is unchanged; the shift of order $N_c/6$ cancels in the virtual part and appears only in the conversion of $\alpha_s$ and $D$. (The sign of the Row~II.IV term was obtained from the general $d_\perp/d$ rule and from this closure, not by re-deriving the row.)

\paragraph{Identification of the $\beta_0$ terms.} The calculation and the result are for pure Yang--Mills theory, $N_f=0$, where $\beta_0/(2N_c)=\frac{11}{6}$ coincides with integrals of the splitting function. Which of the numbers $\frac{11}{6}$ are $\beta_0/(2N_c)$ can be decided by the BLM-type device of adding a quark loop, without computing the $N_f$ terms of the result, which we do not give. With the vertices normalised as here (the paper's triple-gluon matrix element is, exactly, minus the covariant vertex contracted with light-cone polarisation vectors, and the quark vertex is taken in the same normalisation), the angular-averaged squared vertices satisfy $\frac{n_fT_R\sum|V_{q\bar q}|^2}{\frac12N_c\sum|V_{gg}|^2}=\frac{n_fT_R}{N_c}\frac{\xi^2+(1-\xi)^2}{K_{gg}(\xi)}$, $K_{gg}=\xi(1-\xi)-2+\frac{1}{\xi(1-\xi)}=\mathcal P_{gg+}/(2N_c)$, for every $\xi$: a $q\bar q$ loop adds to the kernel $K_{gg}$ of the self-energy the kernel $\frac{n_fT_R}{N_c}[\xi^2+(1-\xi)^2]$ (with its $d_\perp$-dimensional spin trace), and the corresponding $q\bar q$ pair crosses the shock wave in Row~I.V, with colour factor $T_R[U(\v y)-U(\v x)]$ for a point-like pair. Hence the $\frac{11}{3}$ of Row~I.II becomes $\beta_0/N_c$, $\beta_0=\frac{11}{3}N_c-\frac43T_Rn_f$, the $-\frac{11}{6}$ of $N_{2b\,\rm UV}$ becomes $-\beta_0/(2N_c)$, the $\frac{11}{6}$ of the subtraction term in $\mathfrak K^{\rm V}_{2b}$ becomes $\beta_0/(2N_c)$, the $\frac{11}{12}$ of the normalisation, of the coupling counterterm and of the RC term become $\beta_0/(4N_c)$, and the $\frac{\beta_0}{2}$ of $\delta(1-\xi)$ becomes the full $\frac{\beta_0}{2}$: all of these are integrals of the complete one-loop gluon self-energy kernel. In contrast, the coefficient $-\frac{11}{6}$ of the large-distance tail, Eq.~\eqref{eq:dg_tailtot}, is the double-plus integral of the real $\mathcal P_{gg+}$ and stays $-\frac{11}{6}$ for any $N_f$ (the $\delta(1-\xi)$ parts of the DGLAP and RC terms cancel in the tail for any $\beta_0$). The constant of Row~I.II would become $\frac{67}{9}-\frac{\pi^2}{3}-\frac{20}{9}\frac{n_fT_R}{N_c}=\frac{2}{N_c}K_{\rm CMW}$, $K_{\rm CMW}=N_c\big(\frac{67}{18}-\frac{\pi^2}{6}\big)-\frac{10}{9}T_Rn_f$. There is no real $N_f$ term for a measured gluon at this order. The cancellation of the ultraviolet poles between Rows~I.II and I.V, and the pole balance of Table~\ref{tab:poles}, hold for any $\beta_0$.

\paragraph{The collinear term in the convention of CXY.} The two forms of the collinear term, Eq.~\eqref{eq:dg_two}, differ by the term $\Delta$ of Eq.~\eqref{eq:dg_Delta}. The $\delta(1-\xi)$ part of $\Delta$ vanishes, $\Delta$ is integrable at $\xi\to0$ ($\mathcal G_0(\xi)\to0$), and it does not depend on $\mu$. Equivalently $D^{\rm CXY}=\big[\delta(1-\xi)+\frac{\alpha_s}{2\pi}\mathcal P_{gg}(\xi)\log\xi^2\big]\otimes D$. The results of CXY and of~\cite{Altinoluk:2014eka} are meant to be used with standard $\MS$ fragmentation functions, which obey the momentum sum rule and for which the $\MS$ coefficient is the first line of Eq.~\eqref{eq:dg_two}; their form of the collinear term is the coefficient appropriate for the functions $D^{\rm CXY}$, i.e.\ it misses $\Delta$ at $O(\alpha_s)$ when used with standard fragmentation functions. Since $\int_0^1d\xi\,\xi\,\mathcal P_{gg}(\xi)\log\xi^2=4N_c\big(\frac{65}{144}-\zeta_2\big)=\frac{N_c(65-24\pi^2)}{36}\neq0$ ($-14.3225$ for $N_c=3$; with quarks the combination $\int_0^1d\xi\,\xi\log\xi^2[\mathcal P_{gg}+2N_f\mathcal P_{qg}]=\frac{N_c(65-24\pi^2)}{36}-\frac{11N_f}{36}$ would enter), the functions $D^{\rm CXY}$ do not satisfy the momentum sum rule at $O(\alpha_s)$; equivalently, the energy sums of the $e^+e^-$ and $H\to gX$ coefficient functions computed with $\xi^{-2}$ are not those of the total rates. $\Delta$ is given only so that the CXY convention can be recovered; all collinear terms of this paper refer to the first line of Eq.~\eqref{eq:dg_two}. A comparison with the NLO coefficient of hadroproduction, $pp\to gX$, and a collinear (large-$\v k$) limit of the hybrid formula of~\cite{Chirilli:2012jd} have not been carried out; the bare final-state collinear terms of~\cite{Chirilli:2012jd} and of this paper are the same $d$-dimensional integral, so that such a comparison would re-test only the counterterm, which is fixed above.
